\section{Analytic spectral activity, geometry, and antiunitary symmetry}
\label{sec:active}

The preceding theorem fixes a mean clock for an entire family.  This would
be dynamically uninformative if the H\'enon map disappeared under the same
change of variables that proves clock invariance.  Instead, the warp moves
from the potential into the kinetic tensor.  For the centered one-step
nonmagnetic subfamily, we now prove more: every admissible nonzero parameter
\(a\) strictly raises the ground-state energy relative to the radial control
and produces an explicit nonzero relative heat invariant.  The
magnetic parameter acts in a different direction: it leaves the clock
unchanged but modifies the standard geometric time-reversal test.  None of
these facts supplies the missing arithmetic trace.

\subsection{The warp becomes a variable kinetic metric}

Let \(\Psi:\R^2\to\R^2\) be a proper \(C^1\) diffeomorphism with
\(|\det D\Psi|=1\), and define the unitary map
\begin{equation}
 (U_\Psi f)(u)=f(\Psi^{-1}u).
 \label{eq:configuration-unitary}
\end{equation}
Write
\(\HH_{\Psi,0}=-\frac12\Delta+2\pi e^{\pi|\Psi(q)|^2}\).
If \(g=U_\Psi f\), the chain rule and the determinant condition give
\begin{equation}
 \int_{\R^2}|\nabla_q f|^2\dd q
 =\int_{\R^2}\nabla_u g(u)^T G_\Psi(u)\nabla_u g(u)\dd u,
 \label{eq:metric-form}
\end{equation}
where
\begin{equation}
 G_\Psi(u)
 =D\Psi(\Psi^{-1}u)D\Psi(\Psi^{-1}u)^T,
 \qquad \det G_\Psi(u)=1.
 \label{eq:metric-tensor}
\end{equation}
Thus, at \(B=0\),
\begin{equation}
 U_\Psi\HH_{\Psi,0}U_\Psi^{-1}
 =-\frac12\nabla_u\!\cdot G_\Psi(u)\nabla_u
  +2\pi e^{\pi|u|^2}
 \label{eq:unitary-metric}
\end{equation}
in the quadratic-form sense.  The potential is radial in \(u\), but a
nonlinear H\'enon warp leaves a position-dependent determinant-one kinetic
metric.  For the smooth maps used here, if \(G_\Psi\equiv I\), then
\(D\Psi\) is orthogonal everywhere and Euclidean rigidity makes \(\Psi\) an
affine isometry.  Hence the
specific volume-preserving coordinate change does not reduce a nonlinear
warp to the Euclidean radial operator.  The metric identity alone does not
exclude an unrelated accidental isospectrality; equal sublevel volume is
much weaker than equality of spectra.  The next two theorems exclude
isospectrality for the centered one-step nonmagnetic pair, at the bottom of
the spectrum and through a weighted full-spectrum invariant.  No
corresponding all-iterate or magnetic ordering is claimed.  Metric details
are in \cref{app:metric-symmetry}.

\subsection{Strict spectral activity for one centered warp}
\label{sec:strict-ground-state}

To state the analytic activity theorems without changing the normalization
of the full magnetic family, define the auxiliary one-step scalar operators
\begin{equation}
 \mathsf H_{a,\hbar}
 =-\frac{\hbar^2}{2}\Delta+V_{a,1},
 \qquad a>-1,\quad \hbar>0.
 \label{eq:scalar-hbar-family}
\end{equation}
Thus \(\mathsf H_{a,1}=\HH_{a,1,0}\), while
\(\mathsf H_{0,\hbar}\) is the radial equimeasurable control.

\begin{theorem}[Strict one-step ground-state activation]
\label{thm:strict-ground-state}
For every fixed \(a>-1\), \(a\neq0\), and \(\hbar>0\),
\begin{equation}
 \lambda_1(\mathsf H_{a,\hbar})
 >\lambda_1(\mathsf H_{0,\hbar}).
 \label{eq:strict-ground-state}
\end{equation}
Consequently, the two operators are not isospectral, despite having identical
configuration-sublevel measures.  At the paper normalization \(\hbar=1\),
they also have the same two-growing-term Weyl law proved in
\cref{thm:magnetic-weyl}; for general fixed \(\hbar\), both members acquire
the same corresponding semiclassical rescaling.
\end{theorem}

\begin{proof}[Proof architecture]
Because \(\det D\widetilde H_a=1\), the potential \(V_{a,1}\) is
equimeasurable with \(V_{0,1}(q)=2\pi e^{\pi|q|^2}\), and the latter is its
symmetric increasing rearrangement.  If \(f^*\) denotes symmetric decreasing
rearrangement, P\'olya--Szeg\H{o} and the layer-cake form of
Hardy--Littlewood give
\[
 \frac{\hbar^2}{2}\int|\nabla |f|^*|^2
 +\int V_{0,1}(|f|^*)^2
 \le
 \frac{\hbar^2}{2}\int|\nabla f|^2
 +\int V_{a,1}|f|^2
 .
\]
See \citet{LiebLoss2001} for the rearrangement inequalities.  The Rayleigh
principle gives the weak ground-state ordering.  If equality held, the
kinetic and potential rearrangement deficits would be nonnegative with zero
sum, so equality would hold in P\'olya--Szeg\H{o}.  The Brothers--Ziemer
equality classification would then force the positive
analytic warped ground state to be a translate of the radial one
\citep{BrothersZiemer1988}.  Subtracting their eigenvalue equations would
make \(V_{a,1}\) radial about the same point.  This is impossible because
the degree-four homogeneous part of
\[
 \pi^{-1}\log\frac{V_{a,1}(x,y)}{2\pi}
 =x^2+(c_ax+ax^2+y)^2
\]
is \(a^2x^4\), whereas a degree-four polynomial radial about a point has
leading part proportional to \((x^2+y^2)^2\).  The form-domain and equality-
case details are proved in \cref{app:strict-ground-state}.
\end{proof}

At \(a=0\), the centered map is the rotation
\((x,y)\mapsto(-y,x)\), and equality holds.  The theorem orders only the
ground state; it does not order every higher eigenvalue or assert an all-time
heat-trace ordering.  The independent short-time heat statement is proved
next.

\subsection{Uniform relative heat activity}
\label{sec:relative-heat-activity}

The ground-state theorem is a low-energy certificate.  A complementary
weighted full-spectrum certificate appears at high temperature.  Put
\begin{equation}
 \Theta_{a,\hbar}(t)=\Tr(e^{-t\mathsf H_{a,\hbar}}),
 \qquad
 I_a(t)=\int_{\R^2}e^{-tV_{a,1}(q)}
 |\nabla V_{a,1}(q)|^2\dd q,
 \label{eq:heat-activity-definitions}
\end{equation}
and, for \(k=1,2\),
\begin{equation}
 A_k(\lambda)=\int_\lambda^\infty
 we^{-w}\left(\log\frac{w}{\lambda}\right)^k\dd w.
 \label{eq:Ak-definition}
\end{equation}

\begin{theorem}[Relative heat activity of one H\'enon warp]
\label{thm:relative-heat-activity}
Fix \(a>-1\) and \(\hbar>0\), and put
\(L=L(t)=\log(1/(2\pi t))\).  Let \(\gamma\) denote Euler's constant.
As \(t\downarrow0\),
\begin{equation}
 \Theta_{a,\hbar}(t)-\Theta_{0,\hbar}(t)
 =-\frac{t^2}{48\pi}\bigl(I_a(t)-I_0(t)\bigr)
 +O_{a,\hbar}(tL^4),
 \label{eq:relative-heat-carrier}
\end{equation}
where the first-gradient carrier is exactly
\begin{equation}
 I_a(t)-I_0(t)
 =\frac{2a^2}{t^2}
 \left[A_2(2\pi t)+4\pi r_a^2A_1(2\pi t)\right].
 \label{eq:exact-gradient-carrier}
\end{equation}
Consequently,
\begin{align}
 \Theta_{a,\hbar}(t)-\Theta_{0,\hbar}(t)
 =-\frac{a^2}{24\pi}\Bigg[&L^2+
 \bigl(2(1-\gamma)+4\pi r_a^2\bigr)L
 \notag\\
 &+\frac{\pi^2}{6}-2\gamma+\gamma^2
 +4\pi r_a^2(1-\gamma)\Bigg]
 +O_{a,\hbar}(tL^4).
 \label{eq:relative-heat-full-asymptotic}
\end{align}
In particular, for \(a\ne0\),
\begin{equation}
 \lim_{t\downarrow0}
 \frac{\Theta_{a,\hbar}(t)-\Theta_{0,\hbar}(t)}{L^2}
 =-\frac{a^2}{24\pi}<0.
 \label{eq:relative-heat-leading-limit}
\end{equation}
\end{theorem}

For \(a=1.02\), the explicit form is
\begin{equation}
 \Theta_{a,\hbar}(t)-\Theta_{0,\hbar}(t)
 =-0.0137987335661
 \left[L^2+2.98907358486L+1.72992096098\right]
 +O_{\hbar}(tL^4).
 \label{eq:a102-relative-heat}
\end{equation}
This value illustrates the prior-frozen flagship parameter; the theorem
holds for every fixed admissible \(a\) and therefore does not select
\(a=1.02\) arithmetically.

\begin{proof}[Proof architecture]
Write \(Q_a=\widetilde H_a^{-1}\),
\(W(z)=2\pi e^{\pi|z|^2}\), and
\(\varepsilon=\hbar\sqrt t\).  The diagonal Brownian-bridge Feynman--Kac
formula \citep{Simon2005FunctionalIntegration,BoldtGueneysu2023} and the
area-preserving variable \(z=\widetilde H_a(q)\) give
\[
 \Theta_{a,\hbar}(t)=\frac1{2\pi\hbar^2t}
 \int_{\R^2}\mathbb E\exp\!\left[-t\int_0^1
 V_{a,1}(Q_a(z)+\varepsilon\mathbf b_s)\dd s\right]\dd z.
\]
At zero Brownian amplitude, the integrand is \(e^{-tW(z)}\), independently
of \(a\), so the complete classical heat term cancels pointwise.  Bridge
symmetry removes odd amplitude orders; its covariance produces the second-
order coefficient in \cref{eq:relative-heat-carrier}, and angular integration
gives \cref{eq:exact-gradient-carrier}.  For the remainder, split paths at
\(\sup_s|\mathbf b_s|=\delta/(\hbar\sqrt tL)\) and space at
\(\pi|z|^2=L+8\log L\).  Normalized H\'enon derivative estimates bound the
integrated fourth amplitude derivative by \(O_a(L^4)\) on the good event;
Gaussian bridge tails, pathwise Jensen, and incomplete-Gamma decay control
the complement without multiplying a small probability by infinite spatial
volume.  The resulting pre-prefactor remainder is \(O_{a,\hbar}(t^2L^4)\),
which proves \cref{eq:relative-heat-carrier}.  The full argument, including
dominated differentiation on the bad event, is in
\cref{app:relative-heat-proof}.
\end{proof}

The theorem is restricted to one centered warp, zero magnetic field, and
fixed \(\hbar>0\).  The leading coefficient is independent of fixed
\(\hbar\) in two dimensions because the first-gradient factor \(\hbar^2\)
cancels the free heat-kernel prefactor \(\hbar^{-2}\).  The proof does not
cover a varying \(\hbar(t)\), a magnetic Feynman--Kac--It\^o phase, or higher
H\'enon iterates without a new derivative audit.  It proves spectral
activity, not a prime-power trace.

\begin{figure*}[t]
  \centering
  \includegraphics[width=0.96\textwidth]{fig8_relative_heat_asymptotic.pdf}
  \caption{Deterministic consistency checks for the exact relative-heat
  carrier at the prior-fixed value \(a=51/50\).  Write
  \(L=\log(1/(2\pi t))\) and
  \(\mathcal B_a(t)=A_2(2\pi t)+4\pi r_a^2A_1(2\pi t)\).
  Here \(\beta_a=2(1-\gamma)+4\pi r_a^2\) and
  \(\kappa_a=\pi^2/6-2\gamma+\gamma^2+4\pi r_a^2(1-\gamma)\).
  (a) The normalized exact carrier \(\mathcal B_a(t)/L^2\) (markers) is
  compared with \(1+\beta_a/L+\kappa_a/L^2\) (solid curve); the dashed line
  is the limiting value one.  (b) The algebraically stable lower-tail
  representation of the scaled residual approaches the independently
  derived limit
  \(\pi^2(4\pi r_a^2-1)=11.2859411\).  The markers evaluate the exact
  one-dimensional carrier and lower-tail integrals, not the relative heat
  trace.  They neither prove nor measure the \(O_{a,\hbar}(tL^4)\) Brownian
  remainder; that bound is analytic in
  \cref{thm:relative-heat-activity}.}
  \label{fig:relative-heat-carrier}
\end{figure*}

\subsection{The standard geometric time-reversal audit}

Let \(\mathcal C\) denote complex conjugation.  Since the potential and the
chosen vector potential are real,
\begin{equation}
 \mathcal C\HH_{a,n,B}\mathcal C=\HH_{a,n,-B}.
 \label{eq:conjugation-flips-B}
\end{equation}
At \(B=0\), \(\mathcal C\) is an internal antiunitary symmetry and
\(\mathcal C^2=+1\), the standard scalar orthogonal-class structure
\citep{Dyson1962}.  At a fixed \(B\neq0\), bare conjugation is not an
internal symmetry.  An orientation-reversing Euclidean isometry preserving
the potential could still repair it: the reflected conjugation
\(U_S\mathcal C\) also reverses the magnetic pseudoscalar.

\begin{lemma}[No centered reflected-conjugation repair for \(n=1,2\)]
\label{lem:no-reflection}
Let \(a>-1\), \(a\neq0\), and \(n\in\{1,2\}\).  No
orientation-reversing Euclidean isometry \(S\) satisfies
\begin{equation}
 |\widetilde H_a^n(Sq)|^2=|\widetilde H_a^n(q)|^2
 \quad\text{for every }q\in\R^2.
 \label{eq:no-reflection-condition}
\end{equation}
Consequently, at fixed \(B\neq0\), the operator has no time-reversal repair
of the standard reflected-conjugation form \(U_S\mathcal C\).
\end{lemma}

The proof in \cref{app:reflection-proof} uses both the highest homogeneous
term and the quadratic Taylor form at the unique potential minimum.  The
scope is exact: the lemma covers the centered production cases \(n=1,2\),
but it does not rule out every abstract nonlocal antiunitary and does not
claim a result for all \(n\geq3\).  The radial control \(a=0\) retains
reflections, and the uncentered one-iterate H\'enon potential has an
accidental \(x\mapsto-x\) reflection.  Centering therefore matters to this
symmetry audit.  Finally, a scalar magnetic model supplies an
orthogonal-to-unitary test family; it has no spin-\(1/2\) antiunitary with
\(T^2=-1\), so it is not a symplectic/GSE construction.

The signed relative-spectrum container and the prime-power gate are kept
together in \cref{sec:arithmetic-boundary}.  In particular, the heat
invariant proved above is a short-time Laplace-weighted certificate of
nonisospectrality; it is not promoted here to a fixed-time orbit trace or an
arithmetic explicit formula.
